\begin{exsol}{16}{\MR{} Mathematical Exercise~2.1}{53}{Exam}{B1}
\asked Four families share a beach and choose how much to spend on common improvements; find the
efficient total expenditure given each family's value of the spending.

\method Value maximization principle (\MR{} pp.~35--38): with no wealth effects an allocation is
efficient if and only if it maximizes total value, here the sum of the families' values minus the
cost (Q~2.3, item~26, proves this). The improvement is a public good, so the first-order condition
is the Samuelson rule: the \emph{sum} of marginal values equals marginal cost.

\solution
\textbf{Data.} The book's value functions (gross of cost, in dollars) are
\[
  v_1(y) = v_2(y) = 5y - \tfrac12 y^2, \qquad v_3(y) = 7y - \tfrac12 y^2, \qquad v_4(y) = 4y - y^2,
\]
where $y \ge 0$ is total spending; each dollar spent costs a dollar, so $P(y) = -y$ in the
notation of Q~2.3.

\textbf{Total value.} Each family has utility $x_n + v_n(y)$ and the contributions satisfy
$\sum_n x_n = -y$, so total value is
\[
  W(y) \;=\; \sum_{n=1}^{4} v_n(y) - y \;=\; (5+5+7+4)\,y - \big(\tfrac12+\tfrac12+\tfrac12+1\big)\,y^2 - y
  \;=\; 20\,y - \tfrac52\,y^2 .
\]

\textbf{First-order condition.} The marginal values are $v_1' = v_2' = 5 - y$, $v_3' = 7 - y$,
$v_4' = 4 - 2y$. Setting their sum equal to the marginal cost of one dollar:
\[
  \sum_n v_n'(y) = 21 - 5y = 1 \quad\Longrightarrow\quad y^* = 4 .
\]
The second-order condition holds: $W''(y) = -5 < 0$, so $W$ is strictly concave and $y^* = 4$ is the
unique maximizer (and interior, since $W'(0) = 20 > 0$).

\textbf{Values at the optimum.}
\begin{center}
\begin{tabular}{lcccc}
\toprule
Family & 1 & 2 & 3 & 4 \\
\midrule
Gross value $v_n(4)$ & $12$ & $12$ & $20$ & $0$ \\
Marginal value $v_n'(4)$ & $1$ & $1$ & $3$ & $-4$ \\
\bottomrule
\end{tabular}
\end{center}
Total gross value is $44$, the cost is $4$, so $W(4) = 40$. The marginal values sum to $1$, the
marginal cost.

\begin{center}
\begin{tikzpicture}
\begin{axis}[width=0.62\linewidth, height=5.2cm, xmin=0, xmax=5.2, ymin=-7, ymax=22,
  xlabel={total spending $y$}, ylabel={\$ per dollar of $y$}, axis lines=left,
  legend style={font=\footnotesize, draw=none, at={(1,1)}, anchor=north east},
  tick label style={font=\footnotesize}]
  \addplot[thick, hdr, domain=0:5] {21-5*x}; \addlegendentry{$\sum_n v_n'(y) = 21 - 5y$}
  \addplot[thick, ans, domain=0:5] {1}; \addlegendentry{marginal cost $=1$}
  \addplot[dashed, gray, domain=0:5] {4-2*x}; \addlegendentry{family 4: $v_4'(y) = 4 - 2y$}
  \addplot[only marks, mark=*] coordinates {(4,1)};
  \draw[dotted] (axis cs:4,-7) -- (axis cs:4,1);
  \node[font=\footnotesize, above right] at (axis cs:4,1) {$y^* = 4$};
\end{axis}
\end{tikzpicture}
\end{center}

\interp The efficient project is chosen by summing the families' marginal values, not by any one
family's preference: family~4 alone would stop at $y = 2$ (where $v_4' = 0$) and at $y^* = 4$ it
values the improvement at zero, but the other families' high marginal values justify the extra
spending. The value maximization principle says nothing yet about who pays: any contributions
adding up to $4$ are consistent with efficiency, and the division is a separate question (item~17
shows that a bad sharing rule can block the efficient project).

\verify $W'(4) = 20 - 5\cdot 4 = 0$; $W(4) = 80 - 40 = 40$; $W(3.9) = 39.975$ and
$W(4.1) = 39.975$, both below $40$ (\texttt{check\_C.py}).

\answer{Maximize $\sum_n v_n(y) - y = 20y - \tfrac52 y^2$: the sum of marginal values
$21 - 5y$ equals the marginal cost $1$ at $y^* = 4$. Efficient spending is $\$4$ (in the book's
units), creating a total net value of $40$.}
\end{exsol}
